\chapter{Native Async/Await and New Operators (Python 3.5)}

Python 3.5 turned the coroutine pattern of Chapters 9--10 into first-class syntax. \textbf{PEP 492}
gave \texttt{async def} and \texttt{await}---a coroutine that is a \emph{distinct type} from a
generator---along with \texttt{async for} and \texttt{async with}. Two more operators arrived the
same release: \textbf{PEP 465's \texttt{@}} matrix-multiplication operator and \textbf{PEP 448's}
generalized \texttt{*}/\texttt{**} unpacking. This chapter is the canonical home for the
async/await \emph{syntax and object model}; the event-loop internals are Chapter 10 (model) and
Vol IX (depth).

\section*{Section Index}
\begin{itemize}
  \item \textbf{11.1} Native coroutines: \texttt{async def}, \texttt{await}, \texttt{PyCoroObject} (PEP 492)
  \item \textbf{11.2} \texttt{async for} and \texttt{async with}
  \item \textbf{11.3} The matrix-multiplication operator \texttt{@} (PEP 465)
  \item \textbf{11.4} Extended unpacking generalizations (PEP 448)
  \item \textbf{11.5} Performance and anti-patterns
  \item \textbf{11.6} Summary and cross-references
\end{itemize}

\section{Native coroutines: \texttt{async def}, \texttt{await}, and \texttt{PyCoroObject} (PEP 492)}

\textbf{Why this exists.} Through 3.4, a coroutine \emph{was} a generator
(\texttt{@asyncio.coroutine} + \texttt{yield from}). That conflation was dangerous: a coroutine
exposed the iterator protocol (you could \texttt{for}-loop it or \texttt{next()} it outside the
loop), and \texttt{yield from} meant two things. \textbf{PEP 492} split them: \texttt{async def}
produces a coroutine object of a distinct C type (\texttt{PyCoroObject}) that does \emph{not}
implement the iterator protocol, and \texttt{await} is a dedicated keyword.

\begin{lstlisting}[language=Python, caption={async def returns a coroutine --- a distinct type that is not an iterator.}]
import asyncio

async def get_val():
    return 42

c = get_val()                                  # does NOT run the body; returns a coroutine
print("type:", type(c).__name__)
print("has __await__:", hasattr(c, "__await__"), "| has __next__:", hasattr(c, "__next__"))
print("run result:", asyncio.run(get_val()))
c.close()                                       # close the un-awaited coroutine
\end{lstlisting}

Verified output (CPython 3.13.5):

\begin{lstlisting}[caption={Output}]
type: coroutine
has __await__: True | has __next__: False
run result: 42
\end{lstlisting}

\textbf{The \texttt{await} protocol.} \texttt{await expr} requires \texttt{expr} to be
\textbf{awaitable}: a coroutine, or any object whose type implements \texttt{\_\_await\_\_} (the C
\texttt{am\_await} slot) returning an iterator.

\begin{lstlisting}[language=Python, caption={A custom awaitable --- \_\_await\_\_ yields to the loop, then returns a value.}]
import asyncio

class Delayed:
    def __init__(self, value):
        self.value = value
    def __await__(self):
        yield                 # suspend once, handing control to the event loop
        return self.value

async def use_awaitable():
    return await Delayed("custom-result")

print("result:", asyncio.run(use_awaitable()))
\end{lstlisting}

Verified output (CPython 3.13.5):

\begin{lstlisting}[caption={Output}]
result: custom-result
\end{lstlisting}

\textbf{What the interpreter actually does.} \texttt{await} compiles to \texttt{GET\_AWAITABLE}
followed by a \texttt{SEND}/\texttt{YIELD\_VALUE} resume loop---the modern descendant of
\texttt{yield from}.

\begin{lstlisting}[language=Python, caption={Modern await bytecode (3.13) --- GET\_AWAITABLE + SEND/YIELD\_VALUE loop.}]
import dis
async def main():
    val = await get_val()
    return val
dis.dis(main)
\end{lstlisting}

Verified output (CPython 3.13.5, excerpt):

\begin{lstlisting}[caption={Output}]
 14   RETURN_GENERATOR
      POP_TOP
 L1:  RESUME      0
 15   LOAD_GLOBAL 1 (get_val + NULL)
      CALL        0
      GET_AWAITABLE 0
      LOAD_CONST  0 (None)
 L2:  SEND        3 (to L5)       # send into the awaitable; jump to L5 when it finishes
 L3:  YIELD_VALUE 1               # suspend this coroutine, yield to the loop
 L4:  RESUME      3
      JUMP_BACKWARD_NO_INTERRUPT 5 (to L2)
 L5:  END_SEND
      STORE_FAST  0 (val)
\end{lstlisting}

Note \texttt{RETURN\_GENERATOR} at entry (a coroutine is a generator frame under the hood) and the
\texttt{SEND}/\texttt{END\_SEND} pair that replaced the pre-3.11 \texttt{YIELD\_FROM}. Material
showing \texttt{YIELD\_FROM} for \texttt{await} predates 3.11.

\section{\texttt{async for} and \texttt{async with}}

PEP 492 added asynchronous versions for when \emph{advancing an iterator} or \emph{acquiring a
resource} is itself I/O. \texttt{async for} drives \texttt{\_\_aiter\_\_}/\texttt{\_\_anext\_\_}
(a coroutine raising \texttt{StopAsyncIteration}); \texttt{async with} drives
\texttt{\_\_aenter\_\_}/\texttt{\_\_aexit\_\_} (both coroutines). Async comprehensions
(\texttt{[x async for x in ...]}) followed in 3.6 (PEP 530).

\begin{lstlisting}[language=Python, caption={Async iterator and async context manager, consumed by async for / async with.}]
import asyncio

class AsyncRange:
    def __init__(self, n):
        self.n, self.i = n, 0
    def __aiter__(self):
        return self
    async def __anext__(self):
        if self.i >= self.n:
            raise StopAsyncIteration
        self.i += 1
        return self.i - 1

class AsyncCtx:
    async def __aenter__(self):
        print("aenter"); return self
    async def __aexit__(self, *exc):
        print("aexit")

async def demo():
    out = [x async for x in AsyncRange(3)]      # async comprehension (PEP 530, 3.6+)
    async with AsyncCtx():
        pass
    return out

print("collected:", asyncio.run(demo()))
\end{lstlisting}

Verified output (CPython 3.13.5):

\begin{lstlisting}[caption={Output}]
aenter
aexit
collected: [0, 1, 2]
\end{lstlisting}

\section{The matrix-multiplication operator \texttt{@} (PEP 465)}

\textbf{Why this exists.} The numeric community needed a clean infix operator for matrix
multiplication; \texttt{*} was element-wise. \textbf{PEP 465} added \texttt{@} (and \texttt{@=}),
mapping to \texttt{\_\_matmul\_\_}/\texttt{\_\_rmatmul\_\_}/\texttt{\_\_imatmul\_\_} and the C slot
\texttt{nb\_matrix\_multiply}. The operator carries no built-in meaning---it dispatches to the
operands' type.

\begin{lstlisting}[language=Python, caption={@ dispatches to \_\_matmul\_\_; here, a vector dot product.}]
class Vec:
    def __init__(self, data):
        self.data = data
    def __matmul__(self, other):
        return sum(a * b for a, b in zip(self.data, other.data))

print("[1,2,3] @ [4,5,6] =", Vec([1, 2, 3]) @ Vec([4, 5, 6]))

import operator
print("operator.matmul exists:", hasattr(operator, "matmul"))
\end{lstlisting}

Verified output (CPython 3.13.5):

\begin{lstlisting}[caption={Output}]
[1,2,3] @ [4,5,6] = 32
operator.matmul exists: True
\end{lstlisting}

Dispatch follows standard binary-operator rules (try \texttt{\_\_matmul\_\_}; on
\texttt{NotImplemented} try the reflected \texttt{\_\_rmatmul\_\_}; subclass reflected method takes
precedence). NumPy implements \texttt{nb\_matrix\_multiply} to call BLAS/LAPACK routines that
\textbf{release the GIL} (Vol IX), so \texttt{@} on large arrays is readable \emph{and}
parallelizable.

\section{Extended unpacking generalizations (PEP 448)}

\textbf{Why this exists.} \textbf{PEP 448} generalized \texttt{*}/\texttt{**} to multiple
unpackings in collection literals and calls.

\begin{lstlisting}[language=Python, caption={Multiple * and ** unpackings in literals and calls.}]
a, b = [1, 2], [3, 4]
print("[*a, *b, 5]:", [*a, *b, 5])

d1, d2 = {"x": 1}, {"y": 2}
print("{**d1, **d2, 'z': 3}:", {**d1, **d2, "z": 3})

def f(*args):
    return args
print("f(*a, *b):", f(*a, *b))
\end{lstlisting}

Verified output (CPython 3.13.5):

\begin{lstlisting}[caption={Output}]
[*a, *b, 5]: [1, 2, 3, 4, 5]
{**d1, **d2, 'z': 3}: {'x': 1, 'y': 2, 'z': 3}
f(*a, *b): (1, 2, 3, 4)
\end{lstlisting}

\textbf{Modern bytecode.} The 3.5 implementation used \texttt{BUILD\_LIST\_UNPACK}/
\texttt{BUILD\_MAP\_UNPACK}; Python 3.9 replaced those with incremental opcodes
(\texttt{LIST\_EXTEND}, \texttt{LIST\_APPEND}, \texttt{DICT\_UPDATE}, \texttt{DICT\_MERGE}).

\begin{lstlisting}[language=Python, caption={3.9+ builds unpacked literals incrementally.}]
import dis
dis.dis(compile("[*a, *b, 5]", "<s>", "eval"))
\end{lstlisting}

Verified output (CPython 3.13.5):

\begin{lstlisting}[caption={Output}]
  1  BUILD_LIST   0
     LOAD_NAME    0 (a)
     LIST_EXTEND  1
     LOAD_NAME    1 (b)
     LIST_EXTEND  1
     LOAD_CONST   0 (5)
     LIST_APPEND  1
     RETURN_VALUE
\end{lstlisting}

\textbf{Collision rules differ by context.} In a \texttt{dict} literal a duplicate key is silently
overwritten by the rightmost value; in a function \emph{call}, \texttt{f(**d1, **d2)} with
overlapping keys raises \texttt{TypeError}. Dict merge via \texttt{|}/\texttt{|=} (PEP 584) is
Vol V.

\section{Performance and anti-patterns}

\begin{itemize}
  \item \textbf{The un-awaited coroutine}---calling \texttt{async def} without awaiting does
    nothing and warns; always \texttt{await}, schedule, or \texttt{close()} it.
  \item \textbf{Blocking the loop} (Chapter 10)---a sync call inside a coroutine stalls every task.
  \item \textbf{\texttt{@} has no inherent meaning}---it raises \texttt{TypeError} unless the type
    defines \texttt{\_\_matmul\_\_}; it shines with NumPy, not lists.
  \item \textbf{Unpacking copies}---\texttt{[*a, *b]} and \texttt{\{**d1, **d2\}} build new
    objects; avoid in hot loops.
  \item \textbf{\texttt{**} keyword collisions in calls raise}, unlike dict literals.
\end{itemize}

\section{Summary and cross-references}

\begin{itemize}
  \item \textbf{PEP 492} makes \texttt{async def} a \textbf{distinct coroutine type} with
    \texttt{\_\_await\_\_} and no iterator protocol; \texttt{await} compiles to
    \texttt{GET\_AWAITABLE} + \texttt{SEND}/\texttt{YIELD\_VALUE} (the 3.11+ successor to
    \texttt{YIELD\_FROM}).
  \item \textbf{\texttt{async for}/\texttt{async with}} drive the async iterator/context protocols;
    async comprehensions arrived in 3.6.
  \item \textbf{PEP 465} adds \texttt{@}/\texttt{@=}; meaning-free, dispatches to the type---the
    basis of NumPy's BLAS-backed product.
  \item \textbf{PEP 448} generalizes \texttt{*}/\texttt{**} unpacking; 3.9+ compiles it to
    \texttt{LIST\_EXTEND}/\texttt{DICT\_MERGE}.
\end{itemize}

\textbf{Cross-references.} \texttt{yield from} $\rightarrow$ Chapter 9. The event loop, Futures,
Tasks $\rightarrow$ Chapter 10. Internals, \texttt{TaskGroup}, structured concurrency $\rightarrow$
Vol IX. NumPy/BLAS/GIL-releasing numerics $\rightarrow$ Vol IX. \texttt{dict} merge \texttt{|}/
\texttt{|=} $\rightarrow$ Vol V. f-strings and 3.6 ergonomics $\rightarrow$ Vol IV.

\bigskip
\noindent\emph{Volume III complete. Volume IV opens the ``expressive modern Python'' era
(3.6--3.7): f-strings, variable annotations, dataclasses, and context variables.}
